\documentclass[10pt,a4paper]{amsart}
\usepackage[margin=19mm]{geometry}
\usepackage{amsmath,amssymb,mathtools}
\usepackage[scr=rsfs,cal=euler]{mathalfa}
\usepackage{xfrac}
\usepackage[unicode,hidelinks,bookmarks=false]{hyperref}
\newcommand{\ie}{i.e.\ }
\newcommand{\cf}{cf.\ }
\newcommand{\resp}{respectively\ }
\newcommand{\Subsection}{\subsection}
\providecommand{\nobreakdash}{\nobreak\hbox{-}}
\setlength{\emergencystretch}{2em}
\multlinegap0pt
\usepackage[shortlabels]{enumitem}
\usepackage{mathtools}
\usepackage{amssymb}
\usepackage{bm}
\usepackage{xcolor}
\newtheorem{lemma}{Lemma}
\newtheorem{theorem}{Theorem}
\newtheorem{proposition}{Proposition}
\newcommand{\E}{\mathcal{E}}
\newcommand{\F}{\mathcal{F}}
\newcommand{\R}{\mathbb{R}}
\newcommand{\T}{\mathbb{T}}
\newcommand{\Wtwo}{\operatorname{W}_2}
\newcommand{\dotHk}[1][k]{\dot{H}^{#1}}
\newcommand{\potential}[1][\mu]{\dot{H}^1_{#1}}
\newcommand{\rieszmap}[1][\mu]{\mathcal{I}_{#1}}
\title{Feedback Control and Local Convexification of \\ Wasserstein Gradient Flows}
\author{Dante Kalise, Lucas M. Moschen,, Grigorios A. Pavliotis}
\date{Source: arXiv:2603.13588v2 (CC BY 4.0)}
\begin{document}
\maketitle

\noindent\emph{Source and licence.} Feedback Control and Local Convexification of \\ Wasserstein Gradient Flows. Version arXiv:2603.13588v2 (CC BY 4.0), distributed by arXiv under the Creative Commons Attribution 4.0 International licence, http://creativecommons.org/licenses/by/4.0/, as recorded in the arXiv licence metadata for this exact version and shown on the abstract page https://arxiv.org/abs/2603.13588v2. Full licence notice: \texttt{licenses/arxiv-2603.13588-CC-BY-4.0.txt}.\par\smallskip

\noindent\textit{Editorial scope.} Counted unit: the article's proposition prop:constrained\_projected\_hessian (source0.tex lines 1438--1469). The article's complete original proof of it is delivered with it (source0.tex lines 1471--1535, one \texttt{proof} environment of 65 lines). No mathematics is written, repaired or re-derived: the statement and the proof are the article's own lines, in the article's own notation, and no retained line is substituted or re-flowed.  Retained uncounted as the source-local context the proof needs: lines 308--310; lines 312--322; lines 349--363; lines 515--538; lines 563--566; lines 618--624; lines 668--686; lines 1731--1749.  Nothing was omitted from any retained span, so the package carries no omission record.  No retained line contains a citation, so the wrapper carries no bibliography and no bibliography entry.  No retained line contains a figure environment, an image-inclusion command or drawing code, so no image is delivered and the package declares no asset dependency.  Editorial formatting only, in addition: an AMS \texttt{amsart} preamble that restates the article's own declarations and macros used by the retained lines, namely the environments \texttt{lemma}, \texttt{theorem}, \texttt{proposition} on separate counters, together with the standard packages those lines need.  The retained environments print in this document as Lemma 1, Theorem 1, Proposition 1 in order of appearance, whereas the article prints the same items with the numbering of its own full document.  The complete native source of the article as distributed in the arXiv e-print for this exact version is bundled unchanged as \texttt{sources/196342/source0.tex}.
We make the following structural assumptions near $\bar\mu$, which ensure that the Wasserstein Hessian of $\F$ at $\bar\mu$ admits an operator realization on $X$.
\begin{enumerate}[label=(H\arabic*),ref=(H\arabic*)]
    \item\label{ass:H1} The first variation of $\E$ at $\bar\mu$ satisfies $\Phi \in W^{1,\infty}(\T^d)$.

    \item\label{ass:H2} The Wasserstein Hessian of $\E$ at $\bar\mu$, initially defined on $\mathscr{D}$, extends to a bounded symmetric bilinear form on $X$.
    More precisely, for $\varphi, \psi \in W^{2,\infty}(\T^d)$, we set
    \[
    H_\E(\varphi,\psi) \coloneqq \mathrm{Hess}_{\Wtwo} \E(\bar\mu)(\nabla\varphi,\nabla\psi), 
    \]
    and assume that there exists $C_E > 0$ such that for all $\varphi, \psi \in W^{2,\infty}(\T^d)$, 
    \[
    |H_\E(\varphi,\psi)| \le C_E\|\varphi\|_X\|\psi\|_X. 
    \]
    Hence $H_\E$ extends uniquely by continuity to a bounded symmetric bilinear form on $X$.
\end{enumerate}

\begin{lemma}
    \label{lem:mu_bar_bounds}
    Let $\bar\mu \in \operatorname{dom}(\F)$ be a stationary point of $\F$.
    Assume~\ref{ass:H1}.
    Then:
    \begin{itemize}
        \item[(i)] $\bar\mu(x) = Z^{-1}\exp(-\Phi(x)/\sigma)$ for a.e.\ $x\in\T^d$, where $Z \coloneqq \int_{\T^d} e^{-\Phi(y)/\sigma}\,dy$.
        \item[(ii)] There exist constants $0 < c_0 \le C_0 < \infty$ such that for a.e.\ $x \in \T^d$, $c_0 \le \bar\mu(x) \le C_0$.
        \item[(iii)] $\bar\mu\in W^{1,\infty}(\T^d)$ and for a.e.\ $x\in\T^d$,
        \begin{equation}
            \label{eq:derivative_barmu}
            \nabla\bar\mu(x) = -\frac{1}{\sigma} \bar\mu(x) \nabla\Phi(x). 
        \end{equation}
    \end{itemize}
\end{lemma}

\begin{theorem}[Hessian as a closed form]
    \label{thm:hessian_form}
    Let $\bar\mu$ be a stationary point of $\F$ and assume~\ref{ass:H1}--\ref{ass:H2}. 
    Then, there exists a symmetric bilinear form
    \[
    a: D(a) \times D(a) \to \R, \qquad D(a) = \dotHk[2](\T^d),
    \]
    such that, for all $\varphi, \psi \in C^\infty(\T^d)$,
    \begin{equation}
        \label{eq:HessianF_formula}
        \mathrm{Hess}_{\Wtwo} \F(\bar\mu)(\nabla \varphi, \nabla \psi) = a(\varphi,\psi) \coloneqq H_\E(\varphi,\psi) + \sigma q_0(\varphi, \psi).
    \end{equation}
    Moreover, $a$ is densely defined, symmetric, closed, and bounded from below in the Hilbert space $(X, \langle\cdot,\cdot\rangle_X)$.
    Consequently, there exists a unique self-adjoint operator $A : D(A)\subset X \to X$ such that, for all $\varphi \in D(A)$ and $\psi \in D(a)$,
    \begin{equation}
        \label{eq:A_representation}
        a(\varphi, \psi) = \langle A\varphi, \psi\rangle_X.
    \end{equation}
    The domain of $A$ is given by
    \[
    D(A) = \big\{\varphi \in D(a) : \exists f \in X \text{ such that } a(\varphi,\psi) = \langle f,\psi\rangle_X \; \forall \psi \in D(a)\big\},
    \]
    and $A\varphi = f$.
\end{theorem}

\begin{theorem}[Compact resolvent of $A$]
    \label{thm:compact_resolvent_a}
    Under Assumptions~\ref{ass:H1}--\ref{ass:H2}, the self-adjoint operator $A : D(A) \subset X \to X$ has compact resolvent.
\end{theorem}

\begin{enumerate}[label=(H\arabic*),ref=(H\arabic*),start=3]
    \item\label{ass:H3} The map $\xi \mapsto \Theta_E(\Xi(\xi))$ is continuous from $\mathcal{U}_p$ to $X$, and there exists $\mathcal{K}_E \in \mathcal{L}(X', X)$ such that
    \[
    \Theta_E(\Xi(\xi)) = \Theta_E(\bar\mu) + \mathcal{K}_E(\Xi(\xi)-\bar\mu) + o(\|\Xi(\xi)-\bar\mu\|_{X'})
    \]
    in $X$ for $\xi \in \mathcal{U}_p$ sufficiently close to $0$.
\end{enumerate}

\begin{proposition}[Identification of the linearized operator in $X'$]
    \label{prop:linearized_operator_identification}
    Assume~\ref{ass:H1}--\ref{ass:H3}. 
    Then,
    \[
    \left\langle -D\mathcal{G}(\bar\mu)\rieszmap[\bar\mu]\varphi,\psi\right\rangle_{\dot{H}^{-2}, \dotHk[2]} = a(\varphi, \psi)
    \]
    for all $\varphi, \psi\in C^\infty(\T^d)$.
    Moreover, this bilinear form extends by density to all $\varphi, \psi \in D(a)$.
    Set
    \[
    D(L) \coloneqq \left\{\nu \in \rieszmap[\bar\mu]D(a) : \exists f \in X'; \; \forall \psi \in D(a), \; \left\langle D\mathcal{G}(\bar\mu)\nu, \psi\right\rangle_{\dot{H}^{-2}, \dotHk[2]} = \left\langle f,\psi\right\rangle_{X',X} \right\}.
    \]
    For $\nu \in D(L)$, we define $L\nu \coloneqq f$.
    Then
    \[
    D(L) = \rieszmap[\bar\mu]D(A), \qquad L = -\rieszmap[\bar\mu]A\rieszmap[\bar\mu]^{-1}.
    \]
\end{proposition}

\begin{lemma}
    \label{lem:pushforward_map}
    Assume~\ref{ass:H1}.
    For $\xi \in \mathcal{U}_p$, set $\Xi(\xi) = (T_\xi)_{\sharp}\bar\mu$.
    Then $T_\xi$ is a $C^1$-diffeomorphism of $\T^d$.
    Moreover, the translated map $\xi \mapsto \Xi(\xi)-\bar\mu$ maps $\mathcal{U}_p$ into $X'$ and is differentiable in a neighborhood of $0$ and $C^2$ at $0$ as a map into $X'$.
    At $\xi = 0$,
    \[
    D\Xi(0)\eta = -\nabla \cdot (\bar\mu \nabla\eta) = \rieszmap[\bar\mu]\eta,
    \]
    and, for $\eta_1, \eta_2 \in \dotHk[p](\T^d)$ and $\psi \in C^{\infty}(\T^d)$,
    \[
    \left\langle D^2\Xi(0)[\eta_1,\eta_2],\psi\right\rangle_{X',X} = \int_{\T^d}\bar\mu\,\nabla^2\psi\,\nabla\eta_1\cdot\nabla\eta_2\,dx.
    \]
    Moreover, for $\eta_1, \eta_2 \in C^\infty(\T^d)$ and $\zeta \in W^{1,\infty}(\T^d)$,
    \[
    \left\langle D^2\Xi(0)[\eta_1,\eta_2],\zeta\right\rangle_{X',X} = -\sum_{i,j=1}^d\int_{\T^d}\partial_j\zeta\,\partial_i(\bar\mu \,\partial_i\eta_1\,\partial_j\eta_2) \, dx.
    \]
\end{lemma}

\begin{proposition}[Constrained equilibrium and projected linearization]
    \label{prop:constrained_projected_hessian}
    Let $\bar\mu \in \mathcal{M}$ be stationary for $\F$ restricted to $\mathcal{M}$, and assume that Assumptions~\ref{ass:H1}--\ref{ass:H2} hold at $\bar\mu$.
    Define
    \[
    g_\ell^\circ \coloneqq g_\ell - \int_{\T^d} g_\ell \,d\bar\mu, \qquad X_{\mathcal{M}} \coloneqq \left\{\phi \in \potential[\bar\mu] : \langle \phi,g_\ell^\circ\rangle_{\potential[\bar\mu]} = 0, \ell = 1, \ldots, k\right\},
    \]
    and assume that $G^X_{\ell r} \coloneqq \langle g_\ell^\circ, g_r^\circ\rangle_{X}$ is invertible.
    Then the following holds:
    \begin{itemize}
        \item[(i)] there exist unique multipliers $\bar\lambda_1, \ldots, \bar\lambda_k \in \R$ and a constant $c_{\bar\mu} \in \R$ such that
        \begin{equation}
            \label{eq:euler_lagrange_constrained}
            \sigma(\log\bar\mu + 1) + \Theta_E(\bar\mu) + \sum_{\ell=1}^k \bar\lambda_\ell g_\ell = c_{\bar\mu} \quad \text{a.e.\ on } \T^d,
        \end{equation}
        so $\bar\mu \propto \exp \left(-(\Theta_E(\bar\mu) + \sum_{\ell=1}^k \bar\lambda_\ell g_\ell) / \sigma \right)$, and $\bar\mu \in W^{1,\infty}(\T^d)$ satisfies $0 < c_0 \le \bar\mu \le C_0$.
        \item[(ii)] for $\ell = 1, \dots, k$, define 
        \[
        h_\ell(\phi,\psi) \coloneqq \mathrm{Hess}_{\Wtwo} C_\ell(\bar\mu)(\nabla\phi, \nabla\psi) = \int_{\T^d} \nabla^2 g_\ell \nabla \phi \cdot \nabla \psi \, d\bar\mu.
        \]
        The constrained Hessian form
        \[
        a_{\mathcal{M}}(\phi,\psi) \coloneqq a(\phi,\psi) + \sum_{\ell=1}^k \bar\lambda_\ell h_\ell(\phi,\psi), \qquad D(a_{\mathcal{M}}) = D(a) \cap X_{\mathcal{M}},
        \]
        is densely defined, closed, and bounded from below on $X_{\mathcal{M}}$.
        It therefore defines a self-adjoint operator $A_{\mathcal{M}}$ on $X_{\mathcal{M}}$ with compact resolvent.
        \item[(iii)] If, in addition, Assumption~\ref{ass:H3} holds at $\bar\mu$, then the linearized constrained dynamics in $X_{\mathcal{M}}$ admit the weak formulation
        \[
        \langle \partial_t\xi_t, \psi\rangle_X + a_{\mathcal{M}}(\xi_t, \psi) = 0, \qquad \psi \in D(a_{\mathcal{M}}).
        \]
    \end{itemize}
\end{proposition}

\begin{proof}
    For $\phi \in X$, $DC_\ell(\bar\mu)[\rieszmap[\bar\mu]\phi] = \langle\rieszmap[\bar\mu]\phi,g_\ell^\circ\rangle_{X',X} = \langle\phi,g_\ell^\circ\rangle_X$.
    Since $G^X$ is invertible, the differentials $DC_1(\bar\mu), \ldots, DC_k(\bar\mu)$ are linearly independent on $X'$.
    The Lagrange multiplier condition for the restriction of $\F$ to $\mathcal{M}$ therefore gives the existence of $\bar\lambda_1, \ldots, \bar\lambda_k, c_{\bar\mu} \in \R$ such that~\eqref{eq:euler_lagrange_constrained} holds.
    If $(\widetilde\lambda_1, \ldots, \widetilde\lambda_k)$ is another family of multipliers, subtracting both equations and subtracting the average with respect to $\bar\mu$ yields $\sum_{\ell=1}^k(\bar\lambda_\ell - \widetilde\lambda_\ell) g_\ell^\circ = 0$ in $X$.
    Taking the $X$-inner product with $g_r^\circ$ for $r=1,\ldots,k$ gives $G^X(\bar\lambda - \widetilde\lambda) = 0$, and hence $\bar\lambda = \widetilde\lambda$.
    Equation~\eqref{eq:euler_lagrange_constrained} then yields
    \[
    \bar\mu \propto \exp\left(-\frac{1}{\sigma} \left( \Theta_E(\bar\mu)+\sum_{\ell=1}^k\bar\lambda_\ell g_\ell \right) \right).
    \]
    Since $\Theta_E(\bar\mu) \in W^{1,\infty}(\T^d)$ and $g_\ell \in W^{2,\infty}(\T^d)$, (i) follows as in Lemma~\ref{lem:mu_bar_bounds}.

    For $\phi \in C^{\infty}(\T^d)$ and $\mu_s = \Xi(s\phi)$,
    \[
    \frac{d^2}{ds^2} C_\ell(\mu_s)\Big|_{s=0} = \int_{\T^d}\nabla^2g_\ell\,\nabla\phi \cdot \nabla\phi \, d\bar\mu.
    \]
    Polarization gives the expression for $h_\ell$. Moreover, $|h_\ell(\phi,\psi)| \le \|\nabla^2g_\ell\|_{L^\infty} \|\phi\|_X \|\psi\|_X$, so each $h_\ell$ is a bounded symmetric bilinear form on $X$.
    Thus, $a_{\mathcal{M}}$ is a bounded symmetric form perturbation of $a$ restricted to $D(a) \cap X_{\mathcal{M}}$.
    It remains to verify that $D(a_{\mathcal{M}})$ is dense in $X_{\mathcal{M}}$.
    Indeed, the $X$-orthogonal projection is
    \[
    P_{\mathcal{M}}\phi = \phi - \sum_{\ell,r=1}^k ((G^X)^{-1})_{\ell r} \langle \phi,g_r^\circ \rangle_X g_\ell^\circ.
    \]
    Since $g_\ell^\circ \in W^{2,\infty}(\T^d) \subset D(a)$, one has $P_{\mathcal{M}}D(a) \subset D(a) \cap X_{\mathcal{M}}$. 
    The density of $D(a)$ in $X$ therefore implies that $D(a_{\mathcal{M}})$ is dense in $X_{\mathcal{M}}$.
    Since $X_{\mathcal{M}}$ is closed in $X$, the restriction of $a$ to $D(a) \cap X_{\mathcal{M}}$ is closed and bounded from below on $X_{\mathcal{M}}$.
    The same argument as in Theorem~\ref{thm:hessian_form} then shows that $a_{\mathcal{M}}$ is closed and bounded from below.
    The representation theorem therefore defines a self-adjoint operator $A_{\mathcal{M}}$ on $X_{\mathcal{M}}$.
    The embedding $D(a_{\mathcal{M}}) \hookrightarrow X_{\mathcal{M}}$ is compact given that $\dotHk[2](\T^d) \hookrightarrow \dotHk[1](\T^d) \simeq X$ is compact. 
    The proof of Theorem~\ref{thm:compact_resolvent_a} then shows that $A_{\mathcal{M}}$ has compact resolvent.

    Let $\mathcal{G}_{\mathcal{M}}(\mu) \coloneqq \mathcal{G}(\mu) + \sum_{\ell=1}^k m_\ell(\mu)\nabla\cdot(\mu\nabla g_\ell)$, where $m(\mu) = -G(\mu)^{-1}b(\mu)$.
    Assume now~\ref{ass:H3} and define $\E_{\bar\lambda} \coloneqq \E + \sum_{\ell=1}^k \bar\lambda_\ell C_\ell$.
    By parts~(i)--(ii), Assumptions~\ref{ass:H1}--\ref{ass:H2} hold for $\E_{\bar\lambda}$.
    For $\zeta \in \mathcal{U}_p$ ($p > d/2 + 2$), the added constraint terms cancel in first-variation differences, and hence
    \[
    \Theta_{E_{\bar\lambda}}(\Xi(\zeta)) - \Theta_{E_{\bar\lambda}}(\bar\mu) = \Theta_E(\Xi(\zeta)) - \Theta_E(\bar\mu).
    \]
    Thus Assumption~\ref{ass:H3} holds for $\E_{\bar\lambda}$ with the same operator $\mathcal{K}_E$.
    By the Euler--Lagrange equation and integration by parts,
    \[
    b_\ell(\bar\mu) = \int_{\T^d} \bar\mu \nabla g_\ell \cdot \nabla\Theta_E(\bar\mu) \, dx - \sigma\int_{\T^d} \bar\mu \Delta g_\ell \, dx = -\sum_{r=1}^k G^X_{\ell r} \bar\lambda_r.
    \]
    Since $G(\bar\mu) = G^X$, one has $m(\bar\mu) = \bar\lambda$.

    Let $\mu_\varepsilon = \Xi(\varepsilon\phi)$ and $q_\varepsilon = (\mu_\varepsilon - \bar\mu)/\varepsilon$. 
    Assumption~\ref{ass:H3} and Lemma~\ref{lem:pushforward_map} imply $G(\mu_\varepsilon) \to G(\bar\mu)$ and $b(\mu_\varepsilon) \to b(\bar\mu)$ since $\mu_{\varepsilon} \to \bar\mu$ in $L^1$-norm and $(\mu_\varepsilon)_{\varepsilon}$ is uniformly bounded in $L^{\infty}(\T^d)$.
    Since $G(\bar\mu) = G^X$ is invertible, $G(\mu_\varepsilon)$ is invertible for all sufficiently small $\varepsilon$.
    Hence $m(\mu_\varepsilon) \to \bar\lambda$.
    For $\psi \in D(a_{\mathcal{M}})$,
    \[
    \int_{\T^d} \bar\mu \nabla g_\ell \cdot \nabla\psi \, dx = \langle g_\ell^\circ,\psi\rangle_X = 0.
    \]
    It follows that
    \[
    \frac{1}{\varepsilon} \sum_{\ell=1}^k (m_\ell(\mu_\varepsilon) - \bar\lambda_\ell) \left\langle \nabla \cdot (\mu_\varepsilon \nabla g_\ell), \psi \right\rangle_{X',X} = -\sum_{\ell=1}^k (m_\ell(\mu_\varepsilon) - \bar\lambda_\ell) \int_{\T^d}q_\varepsilon \nabla g_\ell \cdot \nabla\psi\,dx \to 0.
    \]
    Therefore, when tested against functions in $D(a_{\mathcal{M}})$, the linearization of $\mathcal{G}_{\mathcal{M}}$ agrees in weak form with that of the gradient flow associated with $\F + \sum_{\ell=1}^k \bar\lambda_\ell C_\ell$.
    Applying Proposition~\ref{prop:linearized_operator_identification} to $\E_{\bar\lambda}$ gives
    \[
    -\left\langle D\mathcal{G}_{\mathcal{M}}(\bar\mu) \rieszmap[\bar\mu]\phi, \psi \right\rangle_{\dot{H}^{-2}, \dotHk[2]} = a_{\mathcal{M}}(\phi, \psi), \qquad \phi \in C^\infty(\T^d) \cap X_{\mathcal{M}}, \psi \in D(a_{\mathcal{M}}).
    \]
    The identity then extends to all $\phi, \psi \in D(a_{\mathcal{M}})$ by density and continuity of the form.
    Therefore, we can write $\langle \partial_t\xi, \psi\rangle_X + a_{\mathcal{M}}(\xi, \psi) = 0$ for $\psi \in D(a_{\mathcal{M}})$, which proves (iii).
\end{proof}

\end{document}
